% Procedencia primaria (lectura de 18-09-2026):
% XROOT = output/EXTREMA_MEDIA_RAZON_NARRACION_Y_REVISION_20260916/ARTICULO/ES/source
% XROOT/libro/01_acoplamiento.tex:51-147,727-749.
% XROOT/sections/k_direccion_dimensional.tex:54-155.
% XROOT/nuclear/12.tex:9-88.
% VIROOT = output/AMPLIACION_FORMAL_LEAN_SERIE_20260916/delivery/USB_ES_EN_STAGE10_20260917/ES/PAQUETES/06_PARTICULAS_Y_MASAS_ES/documentacion_original/payload/source_es
% VIROOT/sections/05_registro_operador_masa.tex:11-119,193-294.
% VIROOT/sections/02_particula_persistencia.tex:210-235.
% XROOT/nuclear/07.tex:5-25,74-154: Hamiltoniano de transporte eliptico.
% XROOT/nuclear/11.tex:878-979: elevacion cotangente y accion discreta.
% Estatutos: cadena K->alpha->accion->masa = RESULTADO_RECUPERADO;
% carta 1+1+10 de eventos y criterio de entrelazamiento = FORMALIZACION_NUEVA;
% congruencia, espectro generalizado y control racional = pruebas completas aqui.

\section{Congruencias positivas, lectores de acción y dinámica hamiltoniana}
\label{sec:k-mass-bridge}

La relación entre el registro dodecafásico y la Ley General de Masas tiene
una composición efectiva. APP produce las dos hojas con sus residuos y
cocientes; TRIT orienta la transición y determina su régimen; TPK compone
selección, transporte y actualización conservando ruta, acarreo, frontera,
supervivencia y memoria. Del estado enriquecido y de su prolongación
compatible proceden conjuntamente las representaciones regionales y el
registro de doce posiciones. La clausura numérica y la lectura material
comparten ese origen y conservan sus operaciones específicas.

Con origen de fase, orientación y base mil declarados, el registro es
\[
 K=(234,543,140,729,659,824,621,58,914,794,146,601).
\]
El lector posicional utiliza
\begin{equation}
 \kappa_{\rm ag}(K)=\frac{\sum_{j=1}^{12}K_j1000^{12-j}}{1000^{12}-1},
 \qquad
 \alpha=\pi_{\rm HMT}+e_{\rm HMT}-\varphi_{\rm HMT}-4-
              \kappa_{\rm ag}(K).
 \label{eq:mass-k-alpha}
\end{equation}
Los valores de la derecha son representaciones internas del mismo proceso
coinductivo. La fórmula expresa el lector de clausura infinita compatible,
con la frontera de acarreo que lo define.

La acción recibe esas representaciones mediante los lectores
\begin{align}
 H_5={}&\frac12\sqrt{1000\alpha/\varphi}-\alpha+
       \frac9{16}\alpha^2-\frac59\alpha^3+
       \frac7{48}\alpha^4-\frac1{54}\alpha^5,\notag\\
 D_A={}&\exp(-100\pi\alpha/9),\qquad
 \mathcal C_\pi=169\alpha^6+\frac{D_A\alpha^7}{1-D_A\alpha},\notag\\
 \eta_{\rm ret}={}&H_5-\frac{90}{\pi}\mathcal C_\pi,
 \qquad R_{\rm act}=\frac{H_5}{\eta_{\rm ret}}>0.
 \label{eq:mass-k-action}
\end{align}
Aquí y en las fórmulas siguientes, \(\pi,\varphi,e\) abrevian las
representaciones HMT ya construidas. La norma determinantal local
\(\Sigma_H=\varphi\alpha^{16}\) fija la década \(-34\), y la
sección \(\hbar_{\rm ret}=\eta_{\rm ret}10^{-34}\mathcal U_S\)
pertenece a la recta de acción declarada. La razón \(R_{\rm act}\)
conserva la relación entre las dos secciones de acción al cambiar su base
dimensional común.

La ruta electrónica central proporciona, mediante sus lectores de
desplazamiento y autocorrelación, el triple \((80,54,6)\). Con
\[
 \Delta_4=\frac{\pi-e\log\pi}{270}\left(1+\frac\pi{729}\right),
 \qquad
 b_e^{\rm alg}=\frac{\sqrt3}{4}
       \bigl[\exp(\varphi/\pi^2)-22\alpha^3\bigr](1+15\Delta_4),
\]
su representación material contiene
\begin{equation}
 \kappa_e=80-54\alpha+6\Delta_4,
 \qquad \varepsilon_e^\beta=b_e^{\rm alg}R_{\rm act}^{\kappa_e}.
 \label{eq:mass-k-electron}
\end{equation}
La composición \eqref{eq:mass-k-alpha}--\eqref{eq:mass-k-electron}
identifica la intervención concreta de \(K\): su lectura de clausura
participa en \(\alpha\), y esa salida actúa después en el lector de
acción y en el exponente de retorno. La elección central del electrón
pertenece a la involución de la coordenatización APP, cuyo único punto fijo es
\((5,5)\). La centralidad de esa celda y su evaluación escalar son
dos resultados enlazados por la ruta y su registro.

\subsection{Coordenadas reversibles del registro de eventos materiales}

El registro material asigna a una ruta \(\gamma\) doce eventos
orientados \(\mathbf e(\gamma)=(e_0,\ldots,e_{11})\), uno por cada
ventana nonádica del calendario de ciento ocho posiciones. El mismo
orden temporal permite realizar estos eventos en el portador de doce
coordenadas utilizado por \(K\). Esta identificación requiere conservar
el origen y la orientación del calendario.

Sean \(\mathbf1\) el vector uniforme,
\(P_{11}=I-\mathbf1\mathbf1^*/12\), y \(P_3\) el proyector
tridimensional de la coordenatización incidencial declarada. La construcción
geométrica del registro determina
\[
 u_K=P_3P_{11}K,\qquad
 \|u_K\|^2=\frac{6638585+2275584\sqrt5}{20}>0,
 \qquad u_K\perp\mathbf1,
\]
y el proyector transversal
\[
 P_{10}(K)=P_{11}-\frac{u_Ku_K^*}{\|u_K\|^2}.
\]

\begin{proposition}[Descomposición material relativa a la dirección de \(K\)]
Para todo vector de eventos \(\mathbf e\), la aplicación
\begin{equation}
 \mathbf e\longmapsto
 \left(s_C=\mathbf1^*\mathbf e,\quad
       a_K=\frac{u_K^*\mathbf e}{\|u_K\|^2},\quad
       v_K=P_{10}(K)\mathbf e\right)
 \label{eq:mass-k-event-chart}
\end{equation}
es invertible sobre su imagen y satisface
\begin{equation}
 \mathbf e=\frac{s_C}{12}\mathbf1+a_Ku_K+v_K,
 \qquad
 \|\mathbf e\|^2=\frac{|s_C|^2}{12}
          +|a_K|^2\|u_K\|^2+\|v_K\|^2.
 \label{eq:mass-k-event-inverse}
\end{equation}
\end{proposition}
\begin{proof}
Los proyectores \(\mathbf1\mathbf1^*/12\),
\(u_Ku_K^*/\|u_K\|^2\) y \(P_{10}(K)\) son ortogonales dos a
dos y suman la identidad. Aplicar esa igualdad a \(\mathbf e\)
demuestra la inversión. La ortogonalidad de los tres sumandos da la
identidad de normas.
\end{proof}

Esta coordenatización conserva el vector de eventos completo. Por ello conserva
también los funcionales que dependen de sus signos, sus sumas y sus
correlaciones, una vez realizada la reconstrucción
\eqref{eq:mass-k-event-inverse}. El registro de predecesores de los nueves,
la cuenta de acción y las coordenadas de torre mantienen su residencia
en la trayectoria enriquecida: la coordenatización anterior modifica únicamente
las coordenadas de los doce eventos. El carácter se evalúa sobre la misma
firma recuperada. La construcción proporciona así una representación
geométrica de la información utilizada por la masa, con una dirección
distinguida por \(K\) y diez componentes transversales.

La cantidad de información que retiene cada representación puede probarse
exactamente. En el dominio algebraico de registros, \(K'=K+\mathbf1\)
satisface
\[
 u_{K'}=u_K,\qquad P_{10}(K')=P_{10}(K),\qquad
 \kappa_{\rm ag}(K')-\kappa_{\rm ag}(K)=\frac1{999}.
\]
La última identidad resulta de sumar la progresión geométrica de los
doce pesos posicionales. Si las representaciones regionales se mantienen
fijas, la lectura de \(\alpha\) cambia en \(-1/999\).
Este control algebraico determina el alcance exacto de la dirección:
su proyector conserva una polarización, mientras la clausura escalar
utiliza además la información uniforme y posicional del registro.
La admisibilidad de \(K'\) como otra historia terminal requiere su
propia construcción APP--TRIT--TPK.

\subsection{Congruencia positiva y espectro del operador material}

Para una multisección \(M\), sea \(\mathscr H_M\) el espacio finito
con base ortonormal etiquetada por sus rutas. El lector de retorno
\(r\in\{D,\Omega\}\) produce un operador autoadjunto
\(B_M^r e_\gamma=\kappa_\gamma^r e_\gamma\). El carácter refinado,
evaluado sobre diferencias respecto del electrón, da el operador basal
\[
 M_M^{(0)}e_\gamma=
 m_e^{(0)}\chi_{\rm ref}(\sigma_\gamma-\sigma_e,
                  \eta_\gamma-\eta_e)e_\gamma.
\]
Los coeficientes de torre y su convergencia pertenecen al lector de cada
ruta. Su representación beta es la congruencia
\begin{equation}
 D_M^r=R_{\rm act}^{B_M^r/2},\qquad
 M_M^{\beta,r}=D_M^rM_M^{(0)}D_M^r.
 \label{eq:mass-k-congruence}
\end{equation}

\begin{proposition}[Invariantes exactos de la congruencia material]
El operador \(D_M^r\) es positivo e invertible. La congruencia
\eqref{eq:mass-k-congruence} conserva rango e inercia; en particular,
conserva la positividad. Si \(M_M^{(0)}\) y \(B_M^r\) son diagonales
en la base de rutas, entonces
\[
 m_\gamma^\beta=m_\gamma^{(0)}R_{\rm act}^{\kappa_\gamma^r}.
\]
Para dos rutas en la misma coordenatización dimensional se obtiene
\[
 \frac{m_Y^\beta}{m_X^\beta}=
 \chi_{\rm ref}(\sigma_Y-\sigma_X,\eta_Y-\eta_X)
 R_{\rm act}^{\kappa_Y^r-\kappa_X^r}.
\]
\end{proposition}
\begin{proof}
El cálculo funcional da
\(D_M^r=\exp((\log R_{\rm act})B_M^r/2)\), con espectro
estrictamente positivo. Su inversa es la exponencial del opuesto.
Para todo \(v\), la forma cuadrática transformada vale
\((D_M^rv)^*M_M^{(0)}(D_M^rv)\). El cambio invertible de variable
preserva las dimensiones máximas de subespacios positivos, negativos
y nulos; éstas determinan la inercia y el rango. En una base común
diagonal, los dos factores aportan conjuntamente
\(R_{\rm act}^{\kappa_\gamma^r}\). Dividir dos valores cancela
la sección dimensional común y utiliza la ley multiplicativa del
carácter.
\end{proof}

La distinción entre transporte de una forma y modificación de un
operador respecto de una métrica fija es comprobable sin aproximaciones.
Para
\[
 M_0=\begin{pmatrix}1&0\\0&4\end{pmatrix},\qquad
 D=\begin{pmatrix}2&0\\0&1/2\end{pmatrix},\qquad
 M'=D^*M_0D=\begin{pmatrix}4&0\\0&1\end{pmatrix},
\]
el determinante de \(D\) es uno y los valores asociados a las dos
rutas se intercambian. Incluso el espectro no ordenado cambia si
\(M_0=I\): se transforma de \(\{1,1\}\) en \(\{4,1/4\}\).
Ambos controles preservan positividad, rango e inercia.

\begin{proposition}[Invarianza espectral con métrica transportada]
Sean \(G_0>0\) una métrica y \(D\) invertible. Si
\[
 M'=D^*M_0D,\qquad G'=D^*G_0D,
\]
los pares \((M',G')\) y \((M_0,G_0)\) tienen el mismo espectro
generalizado, con las mismas multiplicidades.
\end{proposition}
\begin{proof}
Se tiene
\[
 \det(M'-\lambda G')=|\det D|^2\det(M_0-\lambda G_0),
 \qquad (G')^{-1}M'=D^{-1}G_0^{-1}M_0D.
\]
La primera igualdad preserva las raíces y sus multiplicidades; la
segunda exhibe una semejanza de los operadores asociados.
\end{proof}

El flujo diagonal positivo inducido por \(K\) posee, por tanto, dos
realizaciones matemáticas precisas sobre una forma material. Cuando
transporta simultáneamente la forma y su métrica describe el mismo
objeto en otra coordenatización. Cuando actúa sobre la forma conservando la métrica
de rutas, produce una colección operatoria que puede cambiar las masas
relativas. Su identificación física se decide por el lector de rutas y
por la acción declarada.

\subsection{Compatibilidad del flujo dodecafásico con el retorno material}

Sea
\[
 A_K=(\log\rho_A)\operatorname{diag}(\widehat u_{K,1},\ldots,
       \widehat u_{K,12}),\quad
 \widehat u_K=u_K/\|u_K\|,\quad \rho_A=\pi/\varphi^2,
\]
de modo que \(g_K(t)=e^{tA_K}\). La colección de retorno material
en la multisección \(M\) tiene generador
\[
 C_M^r=\tfrac12(\log R_{\rm act})B_M^r,\qquad
 S_M^r(t)=e^{tC_M^r}.
\]

\begin{proposition}[Transporte de la congruencia por un lector entrelazante]
Una aplicación lineal declarada
\(J_M:\mathbb C^{12}\to\mathscr H_M\) que cumple
\begin{equation}
 C_M^rJ_M=J_MA_K
 \label{eq:mass-k-intertwiner}
\end{equation}
satisface
\[
 S_M^r(t)J_M=J_Mg_K(t),
\]
y, para \(M_M(t)=S_M^r(t)M_M^{(0)}S_M^r(t)\),
\begin{equation}
 J_M^*M_M(t)J_M
 =g_K(t)^*\bigl(J_M^*M_M^{(0)}J_M\bigr)g_K(t).
 \label{eq:mass-k-pullback}
\end{equation}
\end{proposition}
\begin{proof}
La identidad \eqref{eq:mass-k-intertwiner} se itera a cada potencia
y se suma en la serie exponencial convergente. Tomar adjuntos y
sustituir en la forma material demuestra
\eqref{eq:mass-k-pullback}.
\end{proof}

El criterio es finito y admite una prueba de esfuerzo directamente
operatoria. En las bases diagonales declaradas, si \(a_j\) y
\(c_\gamma\) son los autovalores de \(A_K\) y \(C_M^r\),
respectivamente, se debe cumplir
\[
 (c_\gamma-a_j)(J_M)_{\gamma j}=0
 \qquad\text{para cada }(\gamma,j).
\]
Cada coeficiente no nulo del lector exige la igualdad correspondiente
de autovalores. La composición recuperada de
\eqref{eq:mass-k-alpha}--\eqref{eq:mass-k-electron} y la coordenatización reversible
\eqref{eq:mass-k-event-chart} son resultados efectivos. La identificación
adicional de sus dos flujos requiere un lector \(J_M\) con la propiedad
\eqref{eq:mass-k-intertwiner}; esta proposición determina exactamente
qué debe comprobar una realización que la sostenga.

\subsection{Transporte elíptico y transformación regular de Legendre}

La realización variacional del transporte tiene un antecedente explícito
en la elipse de acción. El par angular ya expresado determina
\(\eta=\operatorname{artanh}(C^*/A)\), en la cámara
\(A\ne0\), \(|C^*/A|<1\). Las coordenadas \((Q,P)\), ambas de
dimensión raíz de acción, realizan la forma
\[
 \mathcal I_\eta(Q,P)=\frac12\bigl(e^{-\eta}Q^2+e^\eta P^2\bigr),
 \qquad \lambda=P\,dQ,\quad \Omega=dQ\wedge dP.
\]
El transporte entre formas, con
\(M_\eta=\operatorname{diag}(e^{\eta/2},e^{-\eta/2})\) y
\(J_2=\left(\begin{smallmatrix}0&-1\\1&0\end{smallmatrix}\right)\),
es \(M_{\eta'}e^{-\vartheta J_2}M_\eta^{-1}\). Una
interpolación \(\eta(t)\), \(\vartheta(t)\), con
\(\dot\vartheta=\omega(t)>0\), tiene el Hamiltoniano recuperado
del artículo X:
\begin{equation}
 H_{\rm tr}(Q,P,t)=\omega(t)\mathcal I_{\eta(t)}(Q,P)
                 +\frac{\dot\eta(t)}2QP.
 \label{eq:mass-k-hamiltonian}
\end{equation}
El segundo término procede de \(\dot M_\eta M_\eta^{-1}\).
Las ecuaciones
\[
 \dot Q=\omega e^\eta P+\frac{\dot\eta}2Q,
 \qquad
 \dot P=-\omega e^{-\eta}Q-\frac{\dot\eta}2P
\]
conservan \(\mathcal I_{\eta(t)}(Q(t),P(t))\): la variación
explícita de \(\eta\) y la contribución del término de conexión
se cancelan. Este resultado y la identidad
\(P\dot Q-H_{\rm tr}=p\dot q-\omega(q^2+p^2)/2\), para
\(Q=e^{\eta/2}q\), \(P=e^{-\eta/2}p\), están demostrados
en la sección de conservación variacional de la obra de extrema y
media razón holográfica.

\begin{proposition}[Lagrangiano regular del transporte elíptico]
Supongamos \(\eta\in C^2\), \(\omega\in C^1\) y \(\omega>0\).
La transformación de Legendre respecto de \(P\) es invertible y
produce
\begin{equation}
 L_{\rm tr}(Q,\dot Q,t)=
 \frac{\bigl(\dot Q-\dot\eta Q/2\bigr)^2}{2\omega e^\eta}
 -\frac{\omega e^{-\eta}}2Q^2.
 \label{eq:mass-k-lagrangian}
\end{equation}
\end{proposition}
\begin{proof}
Escribamos \(a=\dot\eta/2\), \(b=\omega e^\eta>0\) y
\(c=\omega e^{-\eta}\). Entonces
\(H_{\rm tr}=bP^2/2+aQP+cQ^2/2\), y
\[
 \dot Q=\partial_PH_{\rm tr}=bP+aQ,
 \qquad P=\frac{\dot Q-aQ}{b}.
\]
Al sustituir esta expresión en \(P\dot Q-H_{\rm tr}\), los
términos que contienen \(P\) se reúnen en
\(P(\dot Q-aQ)-bP^2/2=(\dot Q-aQ)^2/(2b)\), lo cual
demuestra \eqref{eq:mass-k-lagrangian}. Además,
\[
 \partial_{\dot Q}L_{\rm tr}=P,
 \qquad
 \partial_P^2H_{\rm tr}=\omega e^\eta>0,
 \qquad
 \partial_{\dot Q}^2L_{\rm tr}=\frac1{\omega e^\eta}>0.
\]
La transformación es regular y sus ecuaciones de Euler--Lagrange
equivalen a las ecuaciones hamiltonianas anteriores.
\end{proof}

Esta eliminación algebraica reúne el Hamiltoniano ya construido con
su forma lagrangiana de segundo orden. La regularidad procede del
término cuadrático de la realización elíptica. Para un levantamiento
cotangente puramente lineal \(H_K(q,p)=p^{\mathsf T}A_Kq\), la
ecuación \(\dot q=A_Kq\) es independiente de \(p\) y el Hessiano
\(\partial_p^2H_K\) es nulo. Su formulación variacional utiliza la
acción de primer orden \(\int p^{\mathsf T}(\dot q-A_Kq)\,dt\).
El corpus contiene también su antecedente discreto
\(L_n=p_{n+1}^{\mathsf T}(q_{n+1}-U_nq_n)\), cuyas ecuaciones
producen el transporte \(q_{n+1}=U_nq_n\) y el transporte dual
\(p_{n+1}=U_n^{-\mathsf T}p_n\). Las dos formulaciones poseen
dominios y operadores determinados. La interpolación temporal y la
sección de acción se conservan como parte de la realización física
declarada.
